%! Unit 2: Quadratic Functions — Practice Test (KEY)
% GENERATED BODY: the blank and key bodies are byte-identical except the header title;
% answers live in \slot / \optok / work, which the preamble shows or hides.
\documentclass[12pt]{article}
\usepackage{algebra2-article}
\usepackage{algebra2-key}   % pulls in -boxes; do NOT also load -boxes
% Coordinate grid: {xmin}{xmax}{ymin}{ymax}
\newcommand{\lgrid}[4]{%
  \draw[step=1, linegray!40, very thin] (#1,#3) grid (#2,#4);
  \draw[->, semithick, charcoal] (#1,0)--(#2,0) node[below right, font=\scriptsize]{$x$};
  \draw[->, semithick, charcoal] (0,#3)--(0,#4) node[left, font=\scriptsize]{$y$};}
% Vertex dot: {x}{y}
\newcommand{\vtx}[2]{\fill[navy] (#1,#2) circle (4.5pt);}
% Part-divider strip (headlinebox takes one arg: the background color)
\newcommand{\parthead}[1]{%
  \boxguard[10]\vspace{6pt}\noindent
  \begin{headlinebox}{forest}\color{white}\bfseries #1\end{headlinebox}%
  \vspace{4pt}}
% THE WORK RULE for a test: \slot{width}{answer} and \opt / \optok are the ONLY
% lines that differ between blank and key, and they differ in the preamble, not the body.
% A slot is a fixed-width underline in both files, so no line rewraps in the key; the
% blank's slot carries a \vphantom of its answer, so a tall answer (a fraction, a root)
% raises the line in both files alike and the page breaks cannot drift.
\newcommand{\slot}[2]{\underline{\makebox[#1]{\ans{#2}}}}
\newcommand{\opt}[2]{(#1)~#2}
% Tighter lead-in above each work block (same in blank and key, so parity holds).
\AtBeginDocument{\setlength{\abovedisplayskip}{4pt}}
\newcommand{\optok}[2]{\mbox{\llap{\textcolor{keyred}{$\boldsymbol{\rightarrow}$}\,}\textcolor{keyred}{\textbf{(#1)}~#2}}}

\begin{document}

\pageheader{Unit 2: Quadratic Functions}{Practice Test --- Answer Key}
\namedateperiod

\vspace{0.06in}
\noindent\textbf{Instructions:} Show all work. Give exact answers --- simplest radical form, and
$a+bi$ form for non-real numbers. No notes or calculator unless your teacher says so.
\hfill\textbf{Total: 100 pts}

\vspace{4pt}
\boxguard[6]
\begin{remindbox}
\small
\textbf{This is a practice test.} It mirrors the real Unit 2 test in length, format, and the
ideas it covers, but the numbers and contexts differ. Work every problem with no notes, then check
your answers against the key.
\end{remindbox}

% ======================================================================
\parthead{Part A --- Reading \& Graphing Parabolas (20 pts)}

\begin{enumerate}[leftmargin=1.9em, itemsep=6pt, label=\arabic*.]
  \item (7 pts) The graph of a quadratic function $f$ is shown. Read each feature from the graph.
    \\[4pt]
    \begin{minipage}[c]{0.40\linewidth}\centering
    \begin{tikzpicture}[scale=0.46]
      \lgrid{-5}{3}{-2}{5}
      \begin{scope}
        \clip (-5,-2) rectangle (3,5);
        \draw[very thick, navy, domain=-4.5:2.5, smooth, samples=80]
          plot (\x, {4-(\x+1)*(\x+1)});
      \end{scope}
      \vtx{-1}{4}
    \end{tikzpicture}
    \end{minipage}%
    \hfill
    \begin{minipage}[c]{0.57\linewidth}
    {\renewcommand{\arraystretch}{1.35}
    \begin{tabular}{@{}r l@{}}
      Vertex: & \slot{2.6cm}{$(-1,\,4)$} \\
      Axis of symmetry: & \slot{2.6cm}{$x=-1$} \\
      Maximum or minimum? & \slot{2.6cm}{maximum} \\
      Its value: & \slot{2.6cm}{$4$} \\
      Zeros: & \slot{3.8cm}{$-3$ and $1$} \\
      Increasing on: & \slot{4.2cm}{$x<-1$, \ $(-\infty,\,-1)$} \\
      Range: & \slot{4.2cm}{$y\le 4$, \ $(-\infty,\,4]$} \\
    \end{tabular}}
    \end{minipage}

  \item (3 pts) Which describes the graph of $y=-2(x-3)^2+5$?
    \\[4pt]
    \opt{A}{opens up; vertex $(3,\,5)$} \hfill \opt{B}{opens down; vertex $(-3,\,5)$}
    \\[4pt]
    \optok{C}{opens down; vertex $(3,\,5)$} \hfill \opt{D}{opens up; vertex $(-3,\,5)$}
\end{enumerate}

\vspace{8pt}
\noindent
\begin{minipage}[t]{0.48\linewidth}
3.\ (5 pts) For $f(x)=2x^2+8x+3$, use $x=-\dfrac{b}{2a}$ to find the axis of symmetry and the
vertex.
\begin{work}
  x &= -\frac{b}{2a} = -\frac{8}{2(2)} = -2 \\
  f(-2) &= 8-16+3 = -5
\end{work}
\vspace{0.12cm}
\hfill Axis: \slot{2.6cm}{$x=-2$}
\\[6pt]
\hspace*{\fill} Vertex: \slot{2.6cm}{$(-2,\,-5)$}
\end{minipage}%
\hfill
\begin{minipage}[t]{0.48\linewidth}
4.\ (5 pts) For $y=(x+4)(x-2)$, find the zeros, the axis of symmetry, and the vertex.
\begin{work}
  x+4 &= 0 \quad\text{or}\quad x-2=0 \\
  x &= \frac{-4+2}{2} = -1 \\
  y &= (-1+4)(-1-2) = -9
\end{work}
\vspace{0.12cm}
\hfill Zeros: \slot{2.6cm}{$-4$ and $2$}
\\[6pt]
\hspace*{\fill} Axis: \slot{2.6cm}{$x=-1$}
\\[6pt]
\hspace*{\fill} Vertex: \slot{2.6cm}{$(-1,\,-9)$}
\end{minipage}

% ======================================================================
\parthead{Part B --- Solving by Factoring (20 pts)}
Factor completely: take out the GCF first. To solve, get the equation equal to zero, factor, and
set each factor equal to zero.

\begin{enumerate}[leftmargin=1.9em, itemsep=6pt, label=\arabic*., start=5]
  \item (8 pts) Factor completely.
    \begin{enumerate}[label=(\alph*), leftmargin=1.9em, itemsep=8pt]
      \item $3x^2-12x-36$
        \begin{work}
          3x^2-12x-36 &= 3(x^2-4x-12)
        \end{work}
        \hfill \slot{4.2cm}{$3(x-6)(x+2)$}
      \item $x^2-81$ \hfill \slot{4.2cm}{$(x+9)(x-9)$}
      \item $x^2-10x+25$ \hfill \slot{4.2cm}{$(x-5)^2$}
      \item $3x^2+4x-4$ \quad (by grouping)
        \begin{work}
          3x^2+4x-4 &= 3x^2+6x-2x-4 \\
                    &= 3x(x+2)-2(x+2)
        \end{work}
        \hfill \slot{4.2cm}{$(3x-2)(x+2)$}
    \end{enumerate}
\end{enumerate}

\vspace{8pt}
\noindent
\begin{minipage}[t]{0.48\linewidth}
6.\ (4 pts) Solve $x^2+3x-28=0$.
\begin{work}
  (x+7)(x-4) &= 0 \\
  x+7 &= 0 \quad\text{or}\quad x-4=0
\end{work}
\vspace{0.12cm}
\hfill Solution: \slot{3.8cm}{$x=-7$ or $x=4$}
\end{minipage}%
\hfill
\begin{minipage}[t]{0.48\linewidth}
7.\ (4 pts) Solve $2x^2=10x$.
\begin{work}
  2x^2-10x &= 0 \\
  2x(x-5) &= 0
\end{work}
\vspace{0.12cm}
\hfill Solution: \slot{3.8cm}{$x=0$ or $x=5$}
\end{minipage}

\begin{enumerate}[leftmargin=1.9em, itemsep=6pt, label=\arabic*., start=8]
  \item (4 pts) Solve $(x+2)(x-3)=6$.
    \begin{work}
      x^2-x-6 &= 6 \\
      x^2-x-12 &= 0 \\
      (x-4)(x+3) &= 0
    \end{work}
    \vspace{0.12cm}
    \hfill Solution: \slot{3.8cm}{$x=4$ or $x=-3$}
\end{enumerate}

% ======================================================================
\parthead{Part C --- Square Roots \& Complex Numbers (20 pts)}
Get the square alone before you take the root, and keep the $\pm$. Write non-real numbers in
$a+bi$ form.

\vspace{6pt}
\noindent
\begin{minipage}[t]{0.48\linewidth}
9.\ (3 pts) Solve $3x^2+5=65$.
\begin{work}
  3x^2 &= 60 \\
  x^2 &= 20 \\
  x &= \pm\sqrt{20}
\end{work}
\vspace{0.12cm}
\hfill Solution: \slot{3.4cm}{$x=\pm2\sqrt{5}$}
\end{minipage}%
\hfill
\begin{minipage}[t]{0.48\linewidth}
10.\ (4 pts) Solve $(x-3)^2=12$.
\begin{work}
  x-3 &= \pm\sqrt{12} \\
  x-3 &= \pm2\sqrt{3}
\end{work}
\vspace{0.12cm}
\hfill Solution: \slot{3.4cm}{$x=3\pm2\sqrt{3}$}
\end{minipage}

\begin{enumerate}[leftmargin=1.9em, itemsep=6pt, label=\arabic*., start=11]
  \item (4 pts) \quad (a) Simplify $\sqrt{-72}$: \slot{2.8cm}{$6i\sqrt{2}$}
    \hfill (b) Solve $x^2+49=0$: \slot{2.8cm}{$x=\pm7i$}

  \item (6 pts) Let $z_1=4+3i$ and $z_2=2-5i$. Write each answer in $a+bi$ form.
    \\[8pt]
    (a) $z_1+z_2=$ \slot{2.8cm}{$6-2i$} \hfill (b) $z_1-z_2=$ \slot{2.8cm}{$2+8i$}
    \\[8pt]
    (c) $z_1\cdot z_2$
    \begin{work}
      (4+3i)(2-5i) &= 8-20i+6i-15i^2 \\
                   &= 8-14i+15
    \end{work}
    \hfill $z_1\cdot z_2=$ \slot{2.8cm}{$23-14i$}

  \item (3 pts) $\sqrt{-4}\cdot\sqrt{-25}=$ \hfill \opt{A}{$10$} \hfill \optok{B}{$-10$} \hfill \opt{C}{$10i$} \hfill \opt{D}{$-10i$}
\end{enumerate}

% ======================================================================
\parthead{Part D --- Completing the Square, Formula \& Discriminant (20 pts)}

\begin{enumerate}[leftmargin=1.9em, itemsep=6pt, label=\arabic*., start=14]
  \item (4 pts) Solve by completing the square: \quad $x^2+8x-5=0$
    \begin{work}
      x^2+8x &= 5 \\
      x^2+8x+16 &= 5+16 \\
      (x+4)^2 &= 21 \\
      x+4 &= \pm\sqrt{21}
    \end{work}
    \vspace{0.12cm}
    \hfill Solution: \slot{3.8cm}{$x=-4\pm\sqrt{21}$}
\end{enumerate}

\vspace{4pt}
\noindent For 15 and 16, use the quadratic formula. Write the equation in standard form first.

\vspace{6pt}
\noindent
\begin{minipage}[t]{0.48\linewidth}
15.\ (5 pts) Solve $2x^2+1=6x$.
\begin{work}
  2x^2-6x+1 &= 0 \\
  x &= \frac{6\pm\sqrt{36-8}}{4} \\
    &= \frac{6\pm\sqrt{28}}{4}
\end{work}
\vspace{0.12cm}
\hfill Solution: \slot{3.4cm}{$x=\tfrac{3\pm\sqrt{7}}{2}$}
\end{minipage}%
\hfill
\begin{minipage}[t]{0.48\linewidth}
16.\ (5 pts) Solve $x^2+6x+13=0$.
\begin{work}
  x &= \frac{-6\pm\sqrt{6^2-4(1)(13)}}{2(1)} \\
    &= \frac{-6\pm\sqrt{-16}}{2} = \frac{-6\pm4i}{2}
\end{work}
\vspace{0.12cm}
\hfill Solution: \slot{3.4cm}{$x=-3\pm2i$}
\end{minipage}

\begin{enumerate}[leftmargin=1.9em, itemsep=6pt, label=\arabic*., start=17]
  \item (6 pts) For each equation, find the discriminant $b^2-4ac$. Then give the number and type
    of solutions (\emph{two real}, \emph{one repeated real}, or \emph{two complex}) and the number
    of $x$-intercepts of its graph.
    \\[6pt]
    {\renewcommand{\arraystretch}{1.9}
    \begin{tabular}{|c|c|c|c|}
      \hline
      \textbf{Equation} & \textbf{Discriminant} & \textbf{Solutions} & \textbf{$x$-intercepts} \\ \hline
      $x^2-6x+9=0$ & \slot{1.6cm}{$0$} & \slot{4.4cm}{one repeated real} & \slot{1.2cm}{$1$} \\ \hline
      $2x^2+3x+4=0$ & \slot{1.6cm}{$-23$} & \slot{4.4cm}{two complex} & \slot{1.2cm}{$0$} \\ \hline
      $x^2+5x-2=0$ & \slot{1.6cm}{$33$} & \slot{4.4cm}{two real} & \slot{1.2cm}{$2$} \\ \hline
    \end{tabular}}
\end{enumerate}

% ======================================================================
\parthead{Part E --- Systems \& Modeling (20 pts)}
Give each solution of a system as an ordered pair. Answer each model question with units.

\vspace{6pt}
\noindent
\begin{minipage}[t]{0.56\linewidth}
18.\ (7 pts) Solve the system by substitution.
\[ y=x^2-4x+1 \qquad y=x-3 \]
\begin{work}
  x^2-4x+1 &= x-3 \\
  x^2-5x+4 &= 0 \\
  (x-1)(x-4) &= 0 \\
  x=1:\quad y &= 1-3 = -2 \\
  x=4:\quad y &= 4-3 = 1
\end{work}
\vspace{0.12cm}
\hfill Solutions: \slot{4.6cm}{$(1,\,-2)$ and $(4,\,1)$}
\end{minipage}%
\hfill
\begin{minipage}[t]{0.40\linewidth}
19.\ (3 pts) How many points do the graphs of $y=x^2+1$ and $y=2x$ have in common?
\\[6pt]
\opt{A}{none} \\[4pt]
\optok{B}{exactly one} \\[4pt]
\opt{C}{exactly two} \\[4pt]
\opt{D}{infinitely many}
\end{minipage}

\begin{enumerate}[leftmargin=1.9em, itemsep=6pt, label=\arabic*., start=20]
  \item (10 pts) At the egg drop, an egg is launched straight up from the roof of the science
    building. Its height, in feet, $t$ seconds after launch is
    \[ h(t)=-16t^2+48t+64. \]
    \begin{enumerate}[label=(\alph*), leftmargin=1.9em, itemsep=8pt]
      \item (2 pts) How high is the roof? \hfill \slot{2.6cm}{$64$ ft}
    \end{enumerate}
\end{enumerate}

\vspace{2pt}
\noindent
\begin{minipage}[t]{0.48\linewidth}
(b)\ (4 pts) When does the egg reach its greatest height, and what is that height?
        \begin{work}
          t &= -\frac{48}{2(-16)} = 1.5 \\
          h(1.5) &= -36+72+64 = 100
        \end{work}
\vspace{0.2cm}
When: \slot{2.6cm}{$1.5$ s} \quad Height: \slot{2.6cm}{$100$ ft}
\end{minipage}%
\hfill
\begin{minipage}[t]{0.48\linewidth}
(c)\ (4 pts) When does the egg hit the ground? Which solution do you reject, and why?
        \begin{work}
          -16t^2+48t+64 &= 0 \\
          t^2-3t-4 &= 0 \\
          (t-4)(t+1) &= 0
        \end{work}
\vspace{0.2cm}
Lands at: \slot{2.6cm}{$t=4$ s}\\[6pt]
Rejected: \slot{4.6cm}{$t=-1$: negative time}
\end{minipage}

\end{document}
